\originalchapter{散射、色散力与腔量子电动力学}
\originalsection{动态极化率与弹性散射}
非共振驱动的二阶能量包含吸收和反旋项. 对$\mathbf E=\mathbf E_0\cos\omega t$, 单方向极化率为
\begin{equation}\alpha_a(\omega)=\frac2\hbar\sum_b\frac{\omega_{ba}|d_{ba}|^2}{\omega_{ba}^2-\omega^2},\qquad U_a=-\frac14\alpha_a(\omega)E_0^2.\label{eq:polar}\end{equation}
靠近某条线且$|\omega-\omega_{ba}|\ll\omega_{ba}$时才可省反旋项. 用$\langle\mathbf d\rangle=\alpha\mathbf E$辐射得到各向同性散射体的总截面$\sigma=\omega^4|\alpha|^2/(6\pi\epsilon_0^2c^4)$. 低频极化率近常数, 因而Rayleigh截面随$\omega^4$增长. 高频利用振子强度和规则得$\alpha\to-e^2/(m\omega^2)$, 回到Thomson截面$8\pi r_e^2/3$, $r_e=e^2/(4\pi\epsilon_0mc^2)$.

\originalsection{van der Waals、共振交换与Casimir}
两个相距$R$的中性原子在非延迟近似有
\[V_{dd}=\frac{\mathbf d_a\cdot\mathbf d_b-3(\mathbf d_a\cdot\hat R)(\mathbf d_b\cdot\hat R)}{4\pi\epsilon_0R^3}.\]
基态无永久偶极, 一阶平均为零. 二阶能量$-\sum_{ij}|\bra{i,j}V_{dd}\ket{g,g}|^2/(\hbar\omega_{ig}+\hbar\omega_{jg})$为负且随$R^{-6}$, 记为$-C_6/R^6$. 对单有效跃迁的各向同性原子,
\[C_6=\frac32\frac{\alpha_a(0)\alpha_b(0)}{(4\pi\epsilon_0)^2}\frac{\hbar\omega_a\omega_b}{\omega_a+\omega_b}.\]
分母是两个虚激发的能量和. 一激发一基态若两原子完全同频, $\ket{eg},\ket{ge}$简并, 一阶对角化就得到$\pm C_3/R^3$; 不同频且$|C_3/R^3|\ll\hbar|\delta|$时, 则回到二阶$R^{-6}$位移. 这就是Förster共振增强的来源.

$R$超过主要跃迁的光传播尺度$c/\omega_a$后, 非延迟静电近似失效, 基态相互作用渐近变成$R^{-7}$的Casimir--Polder形式. 平行理想导体板的单位面积真空能为$-\pi^2\hbar c/(720a^3)$, 压力为$-\pi^2\hbar c/(240a^4)$. 模式和必须相对无板背景做一致正则化; 不能把两个分别发散的零点能直接相减并当作普通有限求和. 真实导体的温度、频率色散与表面条件会改变理想式.

\originalsection{Jaynes--Cummings模型与缀饰态}
单模腔的旋波哈密顿量为$H=\hbar\omega_ca^\dagger a+\hbar\omega_0\sigma_z/2+\hbar g(a\sigma_++a^\dagger\sigma_-)$. 总激发数守恒, 每个$\ket{e,n},\ket{g,n+1}$子空间是耦合$g\sqrt{n+1}$的二能级. 共振分裂为$2\hbar g\sqrt{n+1}$, 真空也有$2\hbar g$分裂; 这不要求初始腔中已有光子.

初态$\ket{e,0}$的激发布居为$\cos^2gt$, 是可逆能量交换. 连续模带来的去相干才使自发发射近似不可逆. 强耦合通常需要$g$大于腔损耗$\kappa$和原子衰减$\Gamma$的相关尺度. 弱耦合坏腔区可绝热消去腔, 产生腔增强发射; 具体Purcell系数还依模式体积、失谐和偏振, 不能仅从“放进腔”就断言增强.

经典强驱动下, 缀饰态间跃迁频率出现$\omega_L$及$\omega_L\pm\Omega_R$, 形成Mollow三重峰. 位置由缀饰能隙给出, 峰权重和线宽须由主方程或两时相关函数求出; 只画三个能级差不能完整确定光谱.

\originalsection{平行板Casimir能的模式和}
8.422的Casimir转写从受边界约束的真空模式出发. 设理想导体板距$a$, 面积趋无限, 横向波矢$k_\parallel$连续, 法向为$n\pi/a$. $n\ge1$有两偏振, $n=0$只一偏振, 所以单位面积零点能可写
\[\frac EA=\frac{\hbar c}{2\pi}\sum_{n=0}^{\infty}{}'
 \int_0^\infty k\sqrt{k^2+(n\pi/a)^2}\dd k.\]
撇号使$n=0$项减半. 每项及总和均发散, 须在模式频率上乘同一个光滑截止$\exp[-\ell\sqrt{k^2+(n\pi/a)^2}]$, 并扣去同体积无板的连续$n$积分. 不能先给每个发散项随意赋值.

记$b=\pi/a$, $F_\ell(n)=\int_0^\infty k\sqrt{k^2+b^2n^2}\exp[-\ell\sqrt{k^2+b^2n^2}]\dd k$. 令$q=\sqrt{k^2+b^2n^2}$, 得$F_\ell(n)=\int_{bn}^\infty q^2\e^{-\ell q}\dd q$. 小$n$展开
\[F_\ell(n)=F_\ell(0)-\frac{b^3n^3}{3}+O(\ell b^4n^4),\]
故$F_\ell'(0)=0$, $F_\ell'''(0)=-2b^3$. Euler--Maclaurin给
\[\sum_{n=0}^\infty{}'F_\ell(n)-\int_0^\infty F_\ell(t)\dd t
 =-\frac{B_4}{4!}F_\ell'''(0)+\cdots
 \longrightarrow-\frac{b^3}{360},\]
用$B_4=-1/30$, 更高奇导数项在$\ell\to0$消失. 因而
\[\frac{E_{\rm Cas}}A=-\frac{\pi^2\hbar c}{720a^3},\qquad
 \frac FA=-\frac{\partial(E_{\rm Cas}/A)}{\partial a}
 =-\frac{\pi^2\hbar c}{240a^4}.\]
负压力说明相吸. 比较模式和与背景的步骤给出了边界怎样改变真空能, 并明确理想导体、零温、无限平板条件; 材料色散和有限温度的修正另需响应函数. 这一板几何结果也不等于直接把原子间$R^{-6}$式延拓到所有距离.

\originalsection{共振荧光: 相干线与Mollow三重峰}
驱动频率$\omega_L$, 共振Rabi频率$\Omega$, 自然布居衰减$\Gamma$, 无额外纯退相位. 光谱来自稳态两时相关
\[S(\nu)=\frac1\pi\Re\int_0^\infty
 \langle\sigma_+(\tau)\sigma_-(0)\rangle_{\rm ss}\e^{-\ii\nu\tau}\dd\tau,
 \qquad \nu=\omega-\omega_L.\]
相关含不衰减常数$|\langle\sigma_-\rangle|^2$, 产生$\nu=0$的相干$\delta$线; 减掉常数后的连接相关给非相干连续谱. 不能把中央窄相干线与中央有限宽荧光峰当作同一贡献.

共振下$H_{\rm rot}=\hbar\Omega X/2$. Bloch方程为
\[\dot u=-\Gamma u/2,\quad
 \dot v=-\Gamma v/2-\Omega w,\quad
 \dot w=\Omega v-\Gamma(w+1).\]
稳态给$\rho_{ee}=\Omega^2/(\Gamma^2+2\Omega^2)$,
$|\langle\sigma_-\rangle|^2=\Omega^2\Gamma^2/(\Gamma^2+2\Omega^2)^2$.
总谱面积为$\rho_{ee}$, 相干比例为$1/(1+s)$, $s=2\Omega^2/\Gamma^2$. 弱驱动主要是弹性相干散射; 强驱动相干比例趋零, 不能说任何强度都只有三条非相干线.

量子回归定理使连接的两时相关按同一齐次Bloch矩阵演化. $u$模的衰减本征值为$-\Gamma/2$, $(v,w)$块的为
\[-3\Gamma/4\pm\ii\Omega',\qquad
 \Omega'=\sqrt{\Omega^2-\Gamma^2/16}.\]
所以强驱动$\Omega\gg\Gamma$时三个非相干峰的中心近$0,\pm\Omega$, 半高半宽为$\Gamma/2,3\Gamma/4,3\Gamma/4$. 令$L_\gamma(\nu)=\gamma/[\pi(\nu^2+\gamma^2)]$, 单位面积归一的领先强驱动谱为
\[\widehat S_{\rm inc}(\nu)\simeq
 \frac12L_{\Gamma/2}(\nu)+
 \frac14L_{3\Gamma/4}(\nu-\Omega)+
 \frac14L_{3\Gamma/4}(\nu+\Omega).\]
中央积分权重是两倍单个边峰, 不是三峰等高. 有限$\Omega/\Gamma$的精确相关还含正弦项及小权重修正; 上式不充作全驱动范围的精确三Lorentz分解.

缀饰态为$(\ket{e,n}\pm\ket{g,n+1})/\sqrt2$, 两支能差约$\hbar\Omega$. 自发跃迁从一激发流形到下一流形, 同支给中央频率, 跨支给两侧频率; 共振强驱动四类跃迁矩阵元的平方均$1/4$, 两条同支线叠在中央, 得上述$2:1:1$面积比. 线宽则来自开放系统的相关衰减, 单算能级差得不到它. 失谐时混合角与支布居改变, 两边峰可不对称.
